\documentclass[11pt,a4paper]{article}
\usepackage[T1]{fontenc}
\usepackage[utf8]{inputenc}
\usepackage{lmodern}
\usepackage{amsmath}
\usepackage[margin=2.3cm]{geometry}
\usepackage{microtype}
\usepackage{booktabs}
\usepackage{tabularx}
\usepackage[font=small,labelfont=bf]{caption}
\usepackage{float}
\usepackage{enumitem}
\usepackage{xcolor}
\usepackage[colorlinks=true,linkcolor=blue!50!black]{hyperref}

\setlist{nosep,leftmargin=*}
\setlength{\parskip}{0.5em}
\setlength{\parindent}{0pt}
\setlength{\emergencystretch}{3em}
\newcolumntype{Y}{>{\raggedright\arraybackslash}X}

\title{\textbf{IRWA Final Project --- Part 2}\\[0.2em]\large Indexing and Evaluation}
\author{\begin{tabular}{ccc}Èric Amargant & Josep Cardona & Jan Balayla\\\small u231867 -- NIA 284566 & \small u231747 -- NIA 284888 & \small u233694 -- NIA 286572\end{tabular}}
\date{November 2026}

\begin{document}
\maketitle

In this part we index the 28,080 products processed in Part~1, rank them with TF-IDF and evaluate the results.
All the Part~2 code is in \texttt{project\_progress/part\_2/}: the notebook \texttt{IRWA\_Part2.ipynb}, the index and
ranking in \texttt{irwa\_indexing.py}, the metrics in \texttt{irwa\_evaluation.py}, and our relevance labels in
\texttt{test\_queries\_labels.csv}.

\textbf{GitHub repository:} \url{https://github.com/Janbacer/Final_Project_IRWA}\\
\textbf{Tag:} \texttt{IRWA-2026-part-2}\\
The \texttt{README.md} in the main folder of the repository explains how to run the code.

\section{Indexing}

\subsection{Inverted index (point 1)}
We do not clean the text again: the notebook loads \texttt{data/processed\_corpus.pkl} from Part~1 and indexes the
field \texttt{all\_terms} (title + product details + sub\_category + brand + description). Queries are cleaned with the
same Part~1 function, \texttt{preprocess\_text}, so query words and product words match.

Each term points to the products that contain it and to its positions in each product:
\texttt{\{term: [[pid, [positions]], \ldots]\}}. Products are identified by their \texttt{pid}.
The index has \textbf{6,343 terms} and 857,225 postings, and it is built in a few seconds.

Queries are \textbf{conjunctive (AND)}: \texttt{search} intersects the posting lists of the query terms, so a product
is returned only if it contains every query term. If one term is not in the index, nothing is returned.

\subsection{Test queries (point 2)}
We looked at the most frequent terms in the \textbf{titles}, because titles are written the way a shopper searches:
gender (\texttt{men}, \texttt{women}), product type (\texttt{tshirt}, \texttt{shirt}, \texttt{pant}, \texttt{jean}),
colour (\texttt{blue}, \texttt{black}, \texttt{white}) and style (\texttt{print}, \texttt{slim}).

Our first five queries used only these words (e.g.\ \emph{men white polo t-shirt}, \emph{women cotton kurta}). Almost
every product in their top~20 was relevant (we labelled 60 of them, all relevant), so simple keyword matching was
enough and the evaluation would not show anything. So each final query \textbf{starts from popular terms and adds one
constraint that keyword matching handles badly}, while staying a realistic query:

\begin{table}[H]
\centering\small
\begin{tabularx}{\textwidth}{@{}c l Y r@{}}
\toprule
\# & query & what makes it hard & matches\\
\midrule
3 & men white formal shirt slim fit & four constraints; \texttt{formal} is often only in the description or the details & 25\\
4 & women floral print shirt & unstitched ``Shirt Fabric'' uses the same words as a shirt & 35\\
5 & women black track pants with zip pockets & zip pockets are only mentioned in the description & 11\\
6 & women pure cotton t-shirt pack of 3 & Part~1 removes numbers, so the ``3'' is lost & 140\\
7 & black running shoes for men & the sub\_category ``Men's Footwear'' also contains women's shoes & 28\\
\bottomrule
\end{tabularx}
\caption{Our test queries and their number of AND matches. They are numbered 3--7 because 1 and 2 are the validation
queries (566 and 225 matches).}
\end{table}

\subsection{TF-IDF ranking (point 3)}
Products and queries are represented with TF-IDF weights: $tf = \text{count} / \lVert D \rVert$ (term count divided by the Euclidean norm of all
the term counts of the product) and $idf = \log(N/df)$. The query gets the same weights, and the score of a product is
the dot product of the query and product vectors. The AND matches are sorted by score; ties are broken by \texttt{pid}.

First observations:
\begin{itemize}
  \item Many products have exactly the same score: in Part~1 we saw that 87\% of the products share their title with
        another one (colour variants, different sellers).
  \item \emph{pack of 3}: packs of 2 and 5 also match, and a pack of 5 ranks first.
  \item \emph{women floral print shirt}: 7 of the top~20 are unstitched shirt fabric.
  \item \emph{black running shoes for men}: 9 of the top~20 are women's shoes.
\end{itemize}

\section{Evaluation}

\subsection{Evaluation metrics (point 1)}
We implemented the following metrics. For one query, with $rel_i = 1$ if the product
at rank $i$ is relevant and $R$ the number of relevant products:
\begin{itemize}
  \item \textbf{P@K}: relevant products in the top $K$, divided by $K$.
  \item \textbf{R@K}: relevant products in the top $K$, divided by $R$.
  \item \textbf{AP@K}: sum of P@$i$ at each rank $i \le K$ that holds a relevant product, divided by the number of relevant products found in the top $K$.
  \item \textbf{F1@K}: harmonic mean of P@K and R@K.
  \item \textbf{MAP}: mean of AP@K over the queries.
  \item \textbf{MRR}: mean over the queries of $1/\text{rank}$ of the first relevant product (0 if none in the top $K$).
  \item \textbf{NDCG@K}: $\sum_{i \le K} (2^{rel_i}-1)/\log_2(i+1)$, divided by the same sum for the ideal ranking.
\end{itemize}

\textbf{How we evaluate.} Only the labelled products can be evaluated, so for each query we keep the labelled products
in the order our engine returns them. The AND search can miss relevant products: they cannot appear in the ranking, but
they still count in $R$ (R@K, F1@K) and in the ideal ranking of NDCG. For query~1, AND returns the 20 labelled products (13 relevant). For query~2 it
returns 8 of them (6 relevant): it misses 4 relevant products, three \emph{Women} jeans without the word \texttt{men}
and one \emph{Tapered Fit} jeans without \texttt{slim}.

\subsection{Results on the validation queries (point 2)}

\begin{table}[H]
\centering\small
\setlength{\tabcolsep}{5pt}
\begin{tabular}{@{}l ccc ccc ccc@{}}
\toprule
 & \multicolumn{3}{c}{$K=5$} & \multicolumn{3}{c}{$K=10$} & \multicolumn{3}{c}{$K=20$}\\
\cmidrule(lr){2-4}\cmidrule(lr){5-7}\cmidrule(l){8-10}
 & Q1 & Q2 & mean & Q1 & Q2 & mean & Q1 & Q2 & mean\\
\midrule
P@K     & 1.000 & 0.800 & 0.900 & 1.000 & 0.600 & 0.800 & 0.650 & 0.300 & 0.475\\
R@K     & 0.385 & 0.400 & 0.392 & 0.769 & 0.600 & 0.685 & 1.000 & 0.600 & 0.800\\
AP@K    & 1.000 & 1.000 & 1.000 & 1.000 & 0.948 & 0.974 & 0.980 & 0.948 & 0.964\\
F1@K    & 0.556 & 0.533 & 0.544 & 0.870 & 0.600 & 0.735 & 0.788 & 0.400 & 0.594\\
RR@K    & 1.000 & 1.000 & 1.000 & 1.000 & 1.000 & 1.000 & 1.000 & 1.000 & 1.000\\
NDCG@K  & 1.000 & 0.869 & 0.934 & 1.000 & 0.716 & 0.858 & 0.995 & 0.716 & 0.855\\
\bottomrule
\end{tabular}
\caption{Validation queries. Mean of AP@K = MAP@K, mean of RR@K = MRR@K.}
\end{table}

\subsection{Our relevance judgments (point 3a)}
We labelled the top~20 results of each of our queries (query~5 has only 11 matches), 91 products in total. A product is
relevant (1) if it is the product type asked, for the gender asked, and meets every attribute of the query. Price,
rating, stock and brand are ignored. Decisions for each query:

\begin{table}[H]
\centering\small
\begin{tabularx}{\textwidth}{@{}c Y r@{}}
\toprule
\# & relevant only if & relevant\\
\midrule
3 & a formal shirt (not casual, not trousers), slim fit, with white among its colours & 5 / 20\\
4 & a women's shirt with a floral print (not shirt fabric, not a T-shirt) & 12 / 20\\
5 & women's black track pants with zip pockets (not men's, not trousers) & 8 / 11\\
6 & exactly a pack of 3 & 6 / 20\\
7 & men's running shoes with black among their colours (not lace-ups, loafers or slip-ons) & 5 / 20\\
\bottomrule
\end{tabularx}
\caption{Labelling decisions and number of relevant products.}
\end{table}

Only the judged products are known, so $R$ is the number of relevant products in the top~20 and R@20 would always be
1. We report $K = 5$ and $10$:

\begin{table}[H]
\centering\small
\setlength{\tabcolsep}{4.5pt}
\begin{tabular}{@{}l cccccc cccccc@{}}
\toprule
 & \multicolumn{6}{c}{$K=5$} & \multicolumn{6}{c}{$K=10$}\\
\cmidrule(lr){2-7}\cmidrule(l){8-13}
query & P & R & AP & F1 & RR & NDCG & P & R & AP & F1 & RR & NDCG\\
\midrule
3    & 0.600 & 0.600 & 0.917 & 0.600 & 1.000 & 0.699 & 0.300 & 0.600 & 0.917 & 0.400 & 1.000 & 0.699\\
4    & 0.800 & 0.333 & 1.000 & 0.471 & 1.000 & 0.869 & 0.900 & 0.750 & 0.928 & 0.818 & 1.000 & 0.915\\
5    & 0.800 & 0.500 & 0.950 & 0.615 & 1.000 & 0.854 & 0.700 & 0.875 & 0.909 & 0.778 & 1.000 & 0.891\\
6    & 0.600 & 0.500 & 0.589 & 0.545 & 0.500 & 0.515 & 0.600 & 1.000 & 0.598 & 0.750 & 0.500 & 0.746\\
7    & 0.200 & 0.200 & 0.500 & 0.200 & 0.500 & 0.214 & 0.300 & 0.600 & 0.341 & 0.400 & 0.500 & 0.414\\
\midrule
mean & 0.600 & 0.427 & 0.791 & 0.486 & 0.800 & 0.630 & 0.560 & 0.765 & 0.739 & 0.629 & 0.800 & 0.733\\
\bottomrule
\end{tabular}
\caption{Our test queries. The mean row gives MAP@K (AP) and MRR@K (RR).}
\end{table}

\subsection{Comments on the metrics (point 3b)}
\begin{itemize}
  \item \textbf{P@K} tells how much of what the user sees is useful. It ignores the order inside the top $K$ and the
        relevant products that were not returned. Query~7 has P@10 = 0.300: most of the top~10 are women's shoes or
        other shoe types.
  \item \textbf{R@K} tells how many of the relevant products we found. It needs to know all the relevant products, so
        it depends on the judgments: for query~2 it never passes 0.600 because AND misses 4 relevant products. For our
        queries it is relative to the judged top~20 only, so it is optimistic.
  \item \textbf{AP@K} rewards relevant products at the top. As defined above, it only averages over the relevant
        products \emph{found}, so it ignores the missed ones: query~2 has AP@10 = 0.948 with R@10 = 0.600,
        and query~3 has AP@10 = 0.917 with P@10 = 0.300 (3 relevant products at ranks 1, 2 and 4, then nothing).
  \item \textbf{F1@K} balances P@K and R@K, so it is low when either of them is low (query~2: 0.600).
  \item \textbf{MAP} and \textbf{MRR} summarise all the queries. MRR only looks at the first relevant product: it is
        1.000 for the validation queries and 0.800 for ours, because queries~6 and~7 start with a non-relevant product
        (a pack of~5, a lace-up shoe).
  \item \textbf{NDCG@K} uses the position of every relevant product and compares with the best possible ranking. It
        could use graded labels; with binary labels it behaves like a position-aware precision. Because the ideal
        ranking includes the missed products, it shows the problem of query~2 (0.716) that AP hides.
\end{itemize}

\textbf{Analysis.} The validation queries look very good (MAP@10 = 0.974), but these numbers only describe the order
\emph{among} the 20 labelled products: in the full ranking of query~1 the labelled products are spread between ranks~6
and~543. Our harder queries give a more realistic picture (MAP@10 = 0.739, mean P@10 = 0.560). Queries~4 and~5 work
well: most of their matches are the product asked. Queries~3, 6 and~7 fail because the words match, but the product
does not meet the need: a casual shirt that mentions ``formal'' in its details, a pack of~5, women's shoes.

\subsection{Limitations and proposed solutions (point 3c)}
\begin{table}[H]
\centering\small
\begin{tabularx}{\textwidth}{@{}>{\raggedright\arraybackslash}p{3.2cm} Y Y@{}}
\toprule
problem & example & possible solution\\
\midrule
AND is too strict (retrieval) & query~2: 4 relevant products are lost because one word is missing (\emph{Women}
jeans, \emph{Tapered Fit} without ``slim'') & rank all products that match most of the terms (OR + ranking, as with
BM25 in Part~3); add synonyms (skinny, tapered $\approx$ slim)\\
All fields count the same (field types) & query~3: a ``Casual Shirt'' matches because its details say ``Model Name:
Formal Shirt''; for query~5, ``zip'' is only in descriptions & index the fields separately and give more weight to the
title and details than to the description\\
Gender and colour words in other fields & query~7: the sub\_category ``Men's Footwear'' also contains women's shoes;
brands like \emph{True Blue} or \emph{Black Beatle} contain colour words & take the gender from the title or
\emph{Ideal For}; use sub\_category and brand as filters, not as text\\
Product type is not checked & query~4: unstitched ``Shirt Fabric'' ranks among the shirts & use the product type of
the details (\emph{Type}, \emph{Generic Name}) as a filter, or require the product word in the title\\
Numbers are removed (query formulation) & query~6: ``pack of 3'' becomes ``pack'', so packs of 2 and 5 match & keep
quantities such as ``pack of 3'' as a structured field and filter on it\\
Ties and duplicates (ranking quality) & many identical titles get the same score; their order is arbitrary & break ties
with rating, discount or stock (Part~3); group colour variants\\
Word counts and text length (ranking quality) & two ``Slim Men Blue Jeans'' score 1.893 and 1.025 only because one
repeats the query words more often & BM25 limits the effect of repeated words and normalises length
better\\
Positions are not used (indexing) & ``full sleeve'' and ``track pants'' are matched as separate words & phrase queries
with the positions already stored in the index\\
\bottomrule
\end{tabularx}
\caption{Main limitations of the current system.}
\end{table}

\section{Assumptions}
\begin{itemize}
  \item The relevance labels of queries~3--7 are our own judgments, following the criteria of Section~2.3.
  \item Products that were not labelled are not used in the evaluation (their relevance is unknown).
  \item Products with the same score are ranked by \texttt{pid}, and the metrics evaluate exactly that order.
  \item AP@K is divided by the relevant products found in the top $K$, not by all the relevant products.
\end{itemize}

\section*{Use of AI tools}
We used AI tools to revise grammar and improve the clarity of our writing, but not to develop or write the core theoretical 
concepts. In coding, we used AI to correct errors and refactor code, always starting from an initial implementation that we 
designed ourselves in the practical sessions. We also used it to test the code (checking the metrics against manual calculations) and to build a small 
web page that let us label the results of our test queries faster. The relevance labels are our own, assigned by hand.


\end{document}
